\documentclass[11pt]{article}
\usepackage[margin=1in]{geometry}
\usepackage{amsmath,amssymb,amsthm,hyperref}
\setlength{\emergencystretch}{2em}
\newtheorem{theorem}{Theorem}
\newtheorem{lemma}{Lemma}
\title{The Euclidean-ball Bessel-integral asymptotic\\MathDB 365832}
\author{Principia Math in Progress}
\date{12 September 2026}
\begin{document}
\maketitle

\paragraph{Task 1: Solve the original problem.}
For fixed \(p>0\), put
\[
 J_n(p)=\int_0^\infty
 \frac{r^{\,n+n/(2p)-1}}{I_{n/2}(r)^{1/p}}\,dr .
\]
Conjecture 5.2 of Mastrantonis \cite[p.~32]{M} asks for its
exponential asymptotic as \(n\to\infty\). The same statement appears as
Conjecture 5.51 in the author's 2025 thesis \cite[p.~214]{T}.
The source proves the case \(p=1\). We prove every fixed \(p>0\).

\begin{theorem}
For every fixed \(p>0\), as the positive integer \(n\) tends to infinity,
\[
 J_n(p)=
 \left[
 n^{1+1/(2p)}e^{-1-1/(2p)}
 \sqrt{p(1+p)^{1+1/p}}
 \right]^n e^{o(n)} .
\]
The error \(o(n)\) here may depend on \(p\).
\end{theorem}

Only the positive Bessel series, its classical Poisson integral,
Stirling's formula, and an elementary Laplace argument are needed.
In particular, the proof controls the unbounded integration interval
separately from the approximation on compact intervals.

\begin{lemma}\label{compact}
As \(\nu\to\infty\),
\[
 \frac{\log I_\nu(\nu z)}{\nu}\longrightarrow
 \eta(z):=\sqrt{1+z^2}+
 \log\frac{z}{1+\sqrt{1+z^2}}
\]
uniformly on each compact subinterval of \(0<z<\infty\).
\end{lemma}
\begin{proof}
The Poisson representation \cite[10.32.2]{DLMF}, for \(\nu>-1/2\), is
\begin{equation}\label{poisson}
 I_\nu(\nu z)=
 \frac{(\nu z/2)^\nu}{\sqrt{\pi}\,\Gamma(\nu+1/2)}
 \int_{-1}^{1}e^{\nu zt}(1-t^2)^{\nu-1/2}\,dt .
\end{equation}
Write the integral as
\(\int_{-1}^1 e^{\nu f_z(t)}w(t)\,dt\), where
\[
 f_z(t)=zt+\log(1-t^2),\qquad
 w(t)=(1-t^2)^{-1/2},\qquad \int_{-1}^1w(t)\,dt=\pi .
\]
The strictly concave function \(f_z\) has its unique maximum at
\[
 t_z=\frac{z}{1+\sqrt{1+z^2}},\qquad
 m(z):=f_z(t_z)
 =\sqrt{1+z^2}-1+\log\frac{2}{1+\sqrt{1+z^2}} .
\]
The integral in \eqref{poisson} is at most \(\pi e^{\nu m(z)}\).
If \(z\) ranges over a fixed compact subinterval of \((0,\infty)\),
the points \(t_z\) remain uniformly separated from \(-1,1\).
For each \(\epsilon>0\), uniform continuity therefore supplies a
fixed \(h>0\) such that
\[
 f_z(t)\ge m(z)-\epsilon\quad
 \text{whenever } |t-t_z|\le h .
\]
Since \(w\ge1\), the integral is at least
\(2h e^{\nu(m(z)-\epsilon)}\).
Its logarithm divided by \(\nu\) consequently tends uniformly to \(m(z)\).
Stirling's formula gives
\[
 \frac1\nu\log
 \frac{(\nu z/2)^\nu}{\sqrt{\pi}\,\Gamma(\nu+1/2)}
 =1+\log(z/2)+o(1),
\]
uniformly on the same compact interval. Adding the two limits proves
the lemma.
\end{proof}

\begin{proof}[Proof of the theorem]
Fix \(p>0\), and abbreviate
\[
 \nu=\frac n2,\qquad a=1+\frac1{2p},\qquad
 K_n=\nu^{-an}J_n(p)
 =\int_0^\infty z^{an-1} I_\nu(\nu z)^{-1/p}\,dz .
\]
On each fixed interval \([\delta,R]\subset(0,\infty)\), Lemma
\ref{compact} gives
\[
 \phi_n(z):=a\log z-\frac{\log I_\nu(\nu z)}{np}
 \longrightarrow
 \phi(z):=a\log z-\frac{\eta(z)}{2p}
\]
uniformly. The weight \(1/z\) is positive and bounded on this interval.
Bounding the integral above by its maximum, and below by integration
over a neighborhood of a maximizer, proves
\begin{equation}\label{laplace}
 \lim_{n\to\infty}\frac1n
 \log\int_\delta^R z^{an-1}I_\nu(\nu z)^{-1/p}\,dz
 =\max_{\delta\le z\le R}\phi(z).
\end{equation}
We next justify that neither omitted tail changes this limit.

\emph{The tail near zero.}
Positivity of the Bessel series implies
\(I_\nu(\nu z)\ge(\nu z/2)^\nu/\Gamma(\nu+1)\).
Since \(an-\nu/p=n\),
\begin{align}
 \int_0^\delta z^{an-1}I_\nu(\nu z)^{-1/p}\,dz
 &\le A_n\frac{\delta^n}{n},\label{zero}\\
 A_n&=\left[\Gamma(\nu+1)(2/\nu)^\nu\right]^{1/p},
 &\lim_{n\to\infty}\frac{\log A_n}{n}
 &=\frac{\log2-1}{2p}.\nonumber
\end{align}
Thus the upper exponential rate of this tail tends to \(-\infty\)
as \(\delta\downarrow0\).

\emph{The tail at infinity.}
For \(n\ge1\), we have \(\nu\ge1/2\).
Restricting \eqref{poisson} to \(t\in[1/2,3/4]\) gives
\[
 I_\nu(\nu z)\ge
 \frac{(\nu z/2)^\nu}{4\sqrt{\pi}\,\Gamma(\nu+1/2)}
 e^{\nu z/2}(7/16)^{\nu-1/2}.
\]
Consequently,
\begin{align}
 \int_R^\infty z^{an-1}I_\nu(\nu z)^{-1/p}\,dz
 &\le B_n\int_R^\infty z^{n-1}e^{-nz/(4p)}\,dz,\label{infty}\\
 B_n&=\left[
 4\sqrt{\pi}\,\Gamma(\nu+1/2)(2/\nu)^\nu
 (7/16)^{-\nu+1/2}\right]^{1/p},\nonumber\\
 \lim_{n\to\infty}\frac{\log B_n}{n}
 &=\frac{\log(32/7)-1}{2p}.\nonumber
\end{align}
For \(R\ge8p\), the inequality
\(\log(z/R)\le(z-R)/R\) yields
\[
 \int_R^\infty z^{n-1}e^{-nz/(4p)}\,dz
 \le\frac{8p}{n}R^{n-1}e^{-nR/(4p)}.
\]
The upper exponential rate in \eqref{infty} is therefore at most
\[
 \frac{\log(32/7)-1}{2p}+\log R-\frac{R}{4p},
\]
which tends to \(-\infty\) as \(R\to\infty\).
These estimates also establish finiteness of \(J_n(p)\) for every
positive integer \(n\).

\emph{The unique maximum and its value.}
Direct differentiation gives
\[
 \eta'(z)=\frac{\sqrt{1+z^2}}{z},\qquad
 \phi'(z)=\frac{2p+1-\sqrt{1+z^2}}{2pz}.
\]
Hence \(\phi\) increases up to
\(z_*=2\sqrt{p(p+1)}\) and decreases thereafter.
Choose \(\delta<z_*<R\) so that both tail rates above are less than
\(\phi(z_*)\).
Equation \eqref{laplace}, the tail estimates, and
\(\log(u+v+w)\le\log3+\max(\log u,\log v,\log w)\)
then show that
\[
 \lim_{n\to\infty}\frac{\log K_n}{n}=\phi(z_*).
\]
Since \(\sqrt{1+z_*^2}=2p+1\), substitution gives
\[
 \phi(z_*)-a\log2
 =-a+\frac12\log p+\frac{1+1/p}{2}\log(1+p).
\]
Finally, \(\nu=n/2\) and \(J_n(p)=\nu^{an}K_n\) yield exactly
the asserted asymptotic.
\end{proof}

\paragraph{Scope and status.}
This solves the full fixed-\(p\) Bessel-integral statement, and hence
also its equivalent Euclidean-ball exponential asymptotic in
\cite[Conjecture 5.1]{M}. It makes no claim to uniformity when \(p\)
depends on \(n\). MathDB was open and its claim fields were empty at the
recorded checks on 12 September 2026. No matching full public claim was
located in the bounded literature audit saved with this proof.

\paragraph{Task 2: Determine the full leading prefactor.}
Replace \(e^{o(n)}\) by an explicit power of \(n\) and an explicit
constant depending on \(p\), with a proved relative error estimate.

\paragraph{Task 3: Analyze varying parameter regimes.}
Determine ranges of sequences \(p=p_n\to0\) or \(p_n\to\infty\)
for which the stated exponential formula remains valid, and identify
any transition asymptotics.

\paragraph{Task 4: Establish a limit law for the dominant radius.}
Normalize the integrand defining \(J_n(p)\) to a probability density.
Prove a central limit theorem about
\(r=n\sqrt{p(p+1)}\), including the variance and a quantitative error.

\paragraph{Task 5: Generalize the order and radial exponent.}
For fixed \(\alpha,\beta,\gamma>0\), determine the asymptotics and
integrability threshold of
\(\int_0^\infty r^{\beta n-1}I_{\alpha n}(r)^{-\gamma}\,dr\),
and analyze the behavior as that threshold is approached.

\begin{thebibliography}{9}
\bibitem{M}
V.~Mastrantonis,
\emph{\(L^p\)-Legendre transforms and Mahler integrals:
Asymptotics and the Fokker--Planck heat flow},
\href{https://arxiv.org/abs/2411.10439}{arXiv:2411.10439v2},
26 November 2024, Conjectures 5.1--5.2.
The author's current publication list describes it as to appear in
\emph{Journal d'Analyse Math\'ematique}.
\bibitem{T}
V.~Mastrantonis,
\emph{From Bergman kernels to Polarity: Perspectives on the
Mahler, Bourgain, and Blocki conjectures},
Ph.D. thesis, University of Maryland, 2025,
\href{https://doi.org/10.13016/cv5w-5vww}{doi:10.13016/cv5w-5vww},
Conjectures 5.50--5.51.
\bibitem{DLMF}
NIST Digital Library of Mathematical Functions,
\href{https://dlmf.nist.gov/10.32.E2}{formula 10.32.2}.
\end{thebibliography}
\end{document}
